\documentclass[a4paper, 12pt]{article}

% Essential packages
\usepackage[T1, T2A]{fontenc}
\usepackage[utf8]{inputenc}
\usepackage[russian, english]{babel}
\usepackage{graphicx}
\usepackage{geometry}
\usepackage{amsmath}
\usepackage{amsfonts}
\usepackage{amssymb}
\usepackage{hyperref}
\usepackage{indentfirst}
\usepackage{authblk}
\usepackage{gensymb}
\usepackage{listings}
\usepackage{xcolor}
\usepackage{minted}
\usepackage{float}
\usepackage{caption}
\usepackage{titlesec}

\captionsetup[figure]{name=Рис., aboveskip=0.5pt, belowskip=0.5pt}


\captionsetup[table]{name=Табл., aboveskip=1pt, belowskip=0.5pt}

\setminted{
    breaklines=true,
    breakautoindent=true,
    breakindent=20pt,
    fontsize=\small,  % опционально: уменьшает шрифт
    frame=lines,      % опционально: рамка для наглядности
    framesep=2mm
}

% \lstset{
%     language=Matlab,
%     basicstyle=\ttfamily\small,
%     numbers=left,
%     numberstyle=\tiny,
%     stepnumber=1,
%     breaklines=true,
%     keywordstyle=\color{blue},
%     commentstyle=\color{green!50!black},
%     stringstyle=\color{red},
%     frame=single
% }

% Page margins setup
\geometry{left=2cm, right=2cm, bottom=2cm, top=2cm}

% Hyperlink setup
\hypersetup{
    colorlinks=true,
    linkcolor=blue,
    filecolor=magenta,
    urlcolor=cyan
}

\titleformat{\section}
  {\newpage\normalfont\Large\bfseries} % Формат заголовка (добавляем \newpage)
  {\thesection.}
  {1em}
  {}  

\begin{document}
\begin{titlepage}

\begin{minipage}[c]{0.65\textwidth}
    {\vspace{0.5cm}\Large\textbf{УНИВЕРСИТЕТ ИТМО}}\\
    \small{Факультет Систем Управления и Робототехники}
\end{minipage}
\begin{minipage}[c]{0.35\textwidth}
    \begin{flushright}
        
\includegraphics[width=0.8\textwidth]{itmo.png}
    \end{flushright}
\end{minipage}

\vfill
\vfill

\begin{center}
\textbf{\Large Лабораторная работа №3 <<Частотные методы>>}\\[5pt]
{\large Частотные методы}


\vfill
\vfill

\begin{flushleft}
    Выполнили:\\
    \text{Бурдасова Виктория: 465306}\\
    \text{Жаворонков Пётр: 465877}\\
    \text{Беляев Александр: 465203}
\end{flushleft}

\vfill

Санкт-Петербург\\
2026

\end{center}
\end{titlepage}
\renewcommand*\contentsname{Содержание}
\tableofcontents

\section{Фильтрация звука}
\subsection{Постановка задачи}
Загрузить с диска файл \texttt{MUHA.wav} и при помощи фильтрации Фурье-образа очистить сигнал от шумов, оставив только голос. 

Для преобразования аудиофайла в массив со значениями сигнала использовалась Python-библиотека librosa. 
Для Фурье-преобразований были использованы методы \texttt{signal.fft: fft, fftshift, ifft, ifftshift.}
Для воспроизведения полученного сигнала как аудиозаписи был использован метод \texttt{sounddevice.play()} (библиотека sounddevice).

\subsection{Ход работы}
Опустим вывод аудиофайла как графика функции сигнала и сразу перейдём к модулю образа. 
Далее под образом будет подразумеваться модуль сдвинутого образа, полученный путём преобразования \texttt{img = fftshift(fft(signal))}.
Образ исходника выглядит следующим образом:
\begin{figure}[H]
    \includegraphics[width=1\textwidth]{pete-graphs/figure_2.pdf}
    \caption{Модуль образа исходного сигнала}
    \label{fig:t2_source_img}
\end{figure}

Перед нами явно была поставлена задача: отделить речь в записи от шумов.
Для выделения речи мы решили обратиться к стандартным в телефонии значениям и положить, что речевая полоса частот лежит в диапазоне от 300 Гц до 3400 Гц.
В связи с тем, что целевой диапазон нам известен, выбор пал на полосовой фильтр: мы отсекаем всё, кроме нужного нам диапазона.

Реализация фильтра в \texttt{Python} выглядит следующим образом:
\begin{minted}{python}
LOW_FREQ = 300    
HIGH_FREQ = 3400
mask = numpy.zeros_like(img_muha_f)
ind = numpy.where((np.abs(freq) < HIGH_FREQ) & (np.abs(freq) > LOW_FREQ))
mask[ind]=1
filtered_muha_img = img_muha_f * mask
\end{minted}

Если посмотреть на модуль отфильтрованного образа (рис. \ref{fig:t2_filtered_img}) сразу видны колоссальные изменения: 
в связи с тем, что многие низкочастотные шумы имели очень высокие амплитуды, после фильтрации осталась только низкоамплитудная информативная часть (максимум уменьшился на 2 порядка).

\begin{figure}[H]
    \includegraphics[width=1\textwidth]{pete-graphs/figure_3.pdf}
    \caption{Модуль образа исходного сигнала}
    \label{fig:t2_filtered_img}
\end{figure}

То же самое видно при сравнении графиков сигнала (рис. \ref{fig:t2_res}) до и после фильтрации.

\begin{figure}[H]
    \includegraphics[width=1\textwidth]{pete-graphs/figure_5.pdf}
    \caption{Модуль образа исходного сигнала}
    \label{fig:t2_res}
\end{figure}

\subsection{Вывод}
При прослушивании получившейся аудиозаписи отчётливо слышно фразу <<Муха -- это маленькая птичка!>>, причем отсутствуют в целом какие-то аудиоартефакты, шумы. Только информация о мухе.

\section{Приложения}

\subsection{Код к заданию 2}
\begin{minted}{python}
from scipy.fft import fft, fftshift, ifft, ifftshift
import matplotlib.pyplot as plt
import numpy as np

from librosa import load
import sounddevice as sd

plt.rcParams.update({'font.size': 16})

def save_separate_figures(extension="pdf", path=""):
    """ Saves each active matplotlib figure to a separate file. """
    # Получаем номера всех созданных фигур
    fig_nums = plt.get_fignums()

    for i in fig_nums:
        plt.figure(i)
        filename = f"figure_{i}.{extension}"
        plt.savefig(path+filename)
        print(f"Сохранен файл: {filename}")

def setup_plot(figsize=(14, 6)):
    plt.figure(figsize=figsize)
    plt.grid()

muha_f, muha_sr = load("lab3/MUHA.wav", mono=True)
dt = 1/muha_sr
T = len(muha_f)*dt
t = np.linspace(0, T, len(muha_f))

freq = fftshift(np.fft.fftfreq(len(muha_f), dt))

setup_plot()
plt.title("Исходный сигнал MUHA.wav")
plt.plot(t, np.real(muha_f), linewidth=2, label="Исходный сигнал")
plt.legend()
plt.xlabel(r"$t$")
plt.ylabel(r"$f(t)$")
img_muha_f = fftshift(fft(muha_f))

setup_plot()
plt.title("Модуль образа исходного сигнала")
plt.plot(freq, np.abs(img_muha_f), linewidth=2, label="Образ исходника")
plt.xlim([-4000, 4000])
plt.legend()
plt.xlabel(r"$t$")
plt.ylabel(r"|$\hat f(t)$|")
HIGH_FREQ = 3400
LOW_FREQ = 300

mask = np.zeros_like(img_muha_f)
ind = np.where((np.abs(freq) < HIGH_FREQ) & (np.abs(freq) > LOW_FREQ))
mask[ind]=1
filtered_muha_img = img_muha_f * mask
mask[ind]=max(np.abs(filtered_muha_img))

setup_plot()
plt.title("Модуль образа после применения фильтра")
plt.plot(freq, np.abs(filtered_muha_img), linewidth=2, label="Фильтрованный образ")
plt.plot(freq, mask, linewidth=1.5, color="orange", label="Полосовой фильтр")
plt.xlim([-4000, 4000])
plt.xlabel(r"$t$")
plt.ylabel(r"|$\hat f(t)$|")
plt.legend()

filtered_muha = ifft(ifftshift(filtered_muha_img))

setup_plot()
plt.title("Отфильтрованный сигнал MUHA.wav")
plt.plot(t, np.real(filtered_muha), linewidth=2, label="Отфильтрованный сигнал")
plt.legend()
plt.xlabel(r"$t$")
plt.ylabel(r"$f(t)$")

setup_plot()
plt.title("Сравнение сигналов до и после обработки")
plt.xlabel(r"$t$")
plt.ylabel(r"$f(t)$, $f_t(t)$")

plt.plot(t, np.real(muha_f), linewidth=4, color="orange", label="Исходный сигнал")
plt.plot(t, np.real(filtered_muha), linewidth=3, color="green", label="Обработанный сигнал")
plt.legend()
plt.fill_between(t, np.real(filtered_muha), np.real(muha_f), alpha=0.4, color="orange")

sd.play(np.real(muha_f), muha_sr)
sd.wait()
sd.play(np.real(filtered_muha), muha_sr)
sd.wait()

save_separate_figures(path="lab3/pete-graphs/")
# plt.show()
\end{minted}

\end{document}
